\documentclass[a5paper,12pt]{extarticle}
\usepackage{array,amsmath,fancybox,amsfonts,graphicx,amsthm,amssymb,wasysym,latexsym,slashed,anysize,subfigure,calc}
\usepackage[spanish]{babel}
\usepackage[utf8]{inputenc}
\usepackage{fancyhdr,float,hyperref,mathtools}
\usepackage{lettrine,multicol,pdflscape,enumerate}

\usepackage[dvipsnames]{xcolor}


%histórico
	\usepackage[owncaptions]{vhistory}
	\renewcommand{\vhhistoryname}{\textcolor{CadetBlue}{Historial de versiones}}
	\renewcommand{\vhversionname}{Versión}
	\renewcommand{\vhdatename}{Fecha}
	\renewcommand{\vhauthorname}{Autor(es)}
	\renewcommand{\vhchangename}{Cambios}
	% Start of the revision history table, after \begin{document}
	%\begin{versionhistory}
 	 %	\vhEntry{1.0}{22.01.04}{JPW|KW}{created}
  	 %	\vhEntry{1.1}{23.01.04}{DP|JPW}{correction}
  	 %  \vhEntry{1.2}{03.02.04}{DP|JPW}{revised after review}
	%\end{versionhistory}	




%ecuaciones tochas
\usepackage{environ}
\NewEnviron{tocho}{%
\begin{equation*}
\scalebox{1.3}{$\BODY$}
\end{equation*}
}

\usepackage{afterpage}
		%uso 1)		%\afterpage{\null\newpage}
		%\newpage
	%uso 2) definiendo un comando y que no cuente la página
		\newcommand\blankpage{%
			\null
			\thispagestyle{empty}%
			\addtocounter{page}{-1}%
			\newpage}



%entorno para química: texto y math mode
%\newcommand*\chem[1]{\textup{#1}}
\newcommand*\chem[1]{\ensuremath{\mathrm{#1}}}

%formatos tachado
\usepackage{ulem}
%    \ sout {text}% tachado
%     \ uwave {text}% línea de onda
%     \ xout {text}% tachado diagonal
%     \ uuline {text}% Doble subrayado


%para los entornos de soluciones
\usepackage{mdframed}
	%uso: 	entorno (begin,end): {mdframed}[backgroundcolor=brown!10] 


%tamaños fuentes extra: 8,9,14,17,20
	\usepackage{extsizes}
	
 	%fuente librebaskerville
 	\usepackage[T1]{fontenc}
	\usepackage{librebaskerville}
	

	
	%fracciones $$ en texto más elaboradas
	 \usepackage{xfrac}
	 
	%permitir que una ecuación larga esté separada por un salto de página
	\allowdisplaybreaks[4]
		
	%gráficos varios 
	\usepackage{graphics}
	
	%fancy tables
	\usepackage{booktabs}
	\usepackage[small,centerlast,up]{caption}

	
	%mejora en la exposición de guarismos
	\usepackage{numprint}
	


\DeclareMathOperator{\lin}{lin}
\DeclareMathOperator{\id}{\textbf{id}}
\DeclareMathOperator{\modulo}{mod}

\newcommand{\mc}[3]{\multicolumn{#1}{#2}{#3}}
\newcommand{\halmos}{\begin{flushright}
\rule{1mm}{2.5mm} \end{flushright}}
\renewcommand{\qed}{\halmos}


%versores: %uso: \uveci\ne\hat{\imath}_{\uvec{i}}+\uvec{j}+\uvec{k}$
\usepackage{bm}
\newcommand{\uveci}{{\bm{\hat{\textnormal{\bfseries\i}}}}}
\newcommand{\uvecj}{{\bm{\hat{\textnormal{\bfseries\j}}}}}
\DeclareRobustCommand{\uvec}[1]{{%
  \ifcsname uvec#1\endcsname
     \csname uvec#1\endcsname
   \else
    \bm{\hat{\mathbf{#1}}}%
   \fi
}}


%algún nuevo operador
\DeclareMathOperator{\cotan}{cotan}


%para textbf en entorno enumerate
\renewcommand{\labelenumi}{%
 \textbf{\theenumi}.
}

%anysize
\marginsize{.8cm}{.8cm}{.7cm}{.5cm}%requiere el paquete anysize,izq,dcha,arriba,abajo


\newtheorem{teo}{{\sc \textbf{Teorema}}}
\newtheorem{defi}{{\sc \textbf{Definición}}}
\newtheorem{coro}{{\sc \textbf{Corolario}}}




\usepackage{everypage,draftwatermark} %Marca de agua

	
%subrayado y highlight. Tiene que estar detrás de usepackage[color]
\usepackage{soul}
	%\textst{normal} 
	\DeclareRobustCommand{\hlcyan}[1]{{\sethlcolor{cyan}\hl{#1}}} %en cyan!
		%\hl{highlighted text} %uso
	

%opening
\title{\textbf{	}}
\date{}
\pagestyle{fancy}
	\fancyhead[LO,RE]{\footnotesize{B2. FÍSICA}---\textcolor{CadetBlue}{v\vhCurrentVersion}}
	\fancyhead[CO,CE]{}
	\fancyhead[RO,LE]{\href{https://classroom.google.com/c/ODc4MDUwMjM1MDY5?cjc=7ke3s4qm}{\nolinkurl{GClassroom: 									7ke3s4qm}}}

	\fancyfoot[LO,RE]{\href{mailto:ricardo.torres@iesjaimeferran.org}{\footnotesize \nolinkurl{ricardo.torres@iesjaimeferran.org}}}
	\fancyfoot[CO,CE]{}



%cajas
\usepackage[skins]{tcolorbox}

\newtcolorbox{tagbox}[1][]{
  colback=blue!10!white,
  arc=0pt,
  auto outer arc,
%  sharp corners,
  boxrule=0pt,
  enhanced jigsaw,
  nobeforeafter,
  width=0.58\textwidth,
  boxsep=0pt,
 % halign=,
  valign=center,
  box align=center,
  colupper={blue!20!black}, % a dark blue
  fontupper={\bfseries},
  top=\fboxsep,
  bottom=\fboxsep,
  left=\fboxsep,
  right=\fboxsep,
  baseline=\fboxsep,
  #1
}

\newtcolorbox{tagboxtwo}[1][]{
  colback=blue!10!white,
  arc=0pt,
  auto outer arc,
%  sharp corners,
  boxrule=0pt,
  enhanced jigsaw,
  nobeforeafter,
  width=.9\textwidth,
  boxsep=0pt,
 % halign=,
  valign=center,
  box align=center,
  colupper={blue!20!black}, % a dark blue
  fontupper={},
  top=\fboxsep,
  bottom=\fboxsep,
  left=\fboxsep,
  right=\fboxsep,
  baseline=\fboxsep,
  #1
}

\newtcolorbox{tagboxtres}[1][]{
  colback=blue!10!white,
  arc=0pt,
  auto outer arc,
%  sharp corners,
  boxrule=0pt,
  enhanced jigsaw,
  nobeforeafter,
  width=.5\textwidth,
  boxsep=0pt,
 % halign=,
  valign=center,
  box align=center,
  colupper={blue!20!black}, % a dark blue
  fontupper={\centering},
  top=\fboxsep,
  bottom=\fboxsep,
  left=\fboxsep,
  right=\fboxsep,
  baseline=\fboxsep,
  #1
}

\newcommand{\sotagbox}[2][]{%
  \begin{tagbox}[#1]%
    #2
  \end{tagbox}%
}

\newcommand{\sotagboxtwo}[2][]{%
  \begin{tagboxtwo}[#1]%
    #2
  \end{tagboxtwo}%
}

\newcommand{\sotagboxtres}[2][]{%
  \begin{tagboxtres}[#1]%
    #2
  \end{tagboxtres}%
}

%uso: \sotagbox[drop shadow]{Modelo nuclear de bola de billar}

%indentado y distancias para párrafos
\usepackage[skip=8pt plus0.5pt, indent=0pt]{parskip}
\addtolength{\skip\footins}{1pc plus 2pt}


\begin{document}
%\maketitle
\pagenumbering{gobble}
\thispagestyle{empty}
\SetWatermarkText{\includegraphics[scale=.23]{yau.png}}
\SetWatermarkScale{1}
\vspace*{\fill}


{\LARGE{B2 FÍSICA}}{\textcolor{CadetBlue}{\vhCurrentVersion}}
\vspace{10pt}

{\large{GUÍA DE LA ASIGNATURA}}
\vspace{1pt}

{\large{Prof. Ricardo Torres Andrés}
}




\newpage
\SetWatermarkLightness{0.96}
\SetWatermarkScale{1}
\SetWatermarkText{\parbox{2\textwidth}{\center Distancia\\B2FIS}}

\blankpage

\tableofcontents

\newpage

\section{¿Qué es y qué no es {\small{FÍSICA}} de {\small{B2}}?}
\vspace{.5cm}

\begin{multicols}{2}
	\hspace*{6mm}\fbox{ \textbf{ES}}
	\vspace*{2mm}
	
		\begin{itemize}
			\item \dots una introducción a la física de cualquier grado/módulo superior científico-técnico y, por ello, es una asignatura \textit{muy práctica}.
			\item \dots una oportunidad para comprender un montón de sucesos de la vida cotidiana relacionada con \textbf{las ondas y la óptica} 
			%(\textit{e.g.} difracción, dispersión, saber qué diablos son las dioptrías de tus gafas\dots),
			 así como con los \textbf{campos} .
			 %(\textit{e.g. }fenómenos de inducción: cargadores inalámbricos, autoinducción\dots ).
			\item \dots el lugar para tener un primer contacto contacto con temas de física del siglo {\footnotesize{XX}} relacionadas con el micromundo y que aparecen de forma casi ubicua en muchas películas, libros y \textit{memes}.
			\item \dots un laboratorio excelente para complementar y aplicar los conocimientos matemáticos de este curso.
			\item \dots una asignatura para  \textbf{ llevar al día}.
		\end{itemize}
	%\vfill\null
	\columnbreak
	%\vfill\null		
	
	
	\begin{flushright}\fbox{ \textbf{NO ES}}\end{flushright}
		\begin{itemize}
			\item \dots un curso riguroso de física: vamos a pasar de puntillas por el formalismo matemático y físico de muchas cosas (aunque los libros se empeñen en \textit{colaros} cosas).
			\item \dots una asignatura a \textit{dejar para el final}. \textbf{Quien avisa no es traidor}.
			\item \dots\textbf{ una asignatura difícil de aprobar}, ni mucho menos imposible. De verdad. 
			En cambio sí es una asignatura que exige una carga importante de trabajo.
		\end{itemize}
\end{multicols}

\vspace{.7cm}

\newpage

\section{Información relevante}
	\subsection{Profesor}

	Mi nombre es \textbf{Ricardo Torres Andrés}, pertenezco al Departamento de Física y Química del {\footnotesize{IES JAIME FERŔAN}} y seré el responsable de la asignatura de {\footnotesize{FÍSICA}} en esta modalidad durante el curso 2026-2027.
	\vspace{1ex}
	
	Además de en las tutorías de la asignatura, puedes contactar conmigo a través de \href{mailto:ricardo.torres@iesjaimeferran.org}{\nolinkurl{ ricardo.torres@iesjaimeferran.org}}	 
{\bfseries\textcolor{red}{a\~nadiendo en el asunto la etiqueta \nolinkurl{[B2-distancia]}}}.
	\subsection{Tutorías y seguimiento telemático}
	
	La asignatura \hl{dispondr\'a de un aula \textit{classroom}} en el que se publicarán materiales relevantes y \hl{se desarrollar\'a a trav\'es de tutor\'ias}. Los estudiantes pueden optar por acudir presencialmente o conectarse al canal \textit{discord} de la asignatura que será publicado en \textit{classroom}.
	\vspace{1ex}
	
	Puedes matricularte con tus credenciales a través del enlace \href{https://classroom.google.com/c/ODA2NTA5NjQxOTA1?cjc=g7p73xbg}{\nolinkurl{g7p73xbg}}. 
	Los datos relevantes son los siguientes:
	\vspace{2ex}
	

	 \textbf{Horario:} M, 15:00-16:30;\hspace{.3cm} \textbf{Lugar:} Aula 1.4;\hspace{.3cm}\href{https://iesjaimeferran.org/wp-content/uploads/calendario-Distancia-26-27.pdf}{\texttt{Calendario del curso}}


	\subsection{Normativa}
Las competencias, criterios de calificación y contenidos son los recogidos en el \href{https://www.bocm.es/boletin/CM_Orden_BOCM/2022/07/26/BOCM-20220726-1.PDF}{\texttt{Decreto 65/2022, de 20 de julio de 2022}}, del Consejo de Gobierno, por el que se establece para la Comunidad de Madrid el currículo del Bachillerato.

	
	\subsection{Sobre la bibliografía}
 \hl{No se propone libro de texto}: la asignatura tendrá una componente teórica basada en los contenidos que se desarrollen en las tutorías así como a través de los materiales provistos en el \textit{classroom} de la asignatura. 
 
 
 La parte práctica consistirá \textit{fundamentalmente} en ejercicios de {\footnotesize PAU} y {\footnotesize EvAU} clasificados en  \textit{dossiers} para cada tema. Enrique García mantiene un repositorio actualizado de problemas de los últimos años que puede ser útil: puedes encontrarlo en la célebre \href{https://www.fiquipedia.es/home/pruebas/pruebasaccesouniversidad/paufisica/}{\nolinkurl{FIQUIPEDIA}}.


Pese a esto, y debido a la naturaleza del curso que nos ocupa y el número de horas de trabajo individual que se espera del estudiante, \hl{nunca viene mal tener a mano un manual del nivel del curso (\textit{cualquier texto de cualquier editorial valdr\'a})} para tener un repositorio de ejercicios resueltos con suficiente detalle. Mi consejo es que te hagas con uno de estos libros (pídelo prestado o cómpralo barato de segunda mano).


\section{Método de trabajo}
\hl{Las tutorías \textbf{no son clases de teor\'ia}, sino sesiones de gu\'ia al estudio y atenci\'on a dudas}. Basta comparar las horas semanales del curso ordinario ($4\,h$) con las de la modalidad {\footnotesize DISTANCIA} ($1.5\,h$) para entender el porqué. Considerando que el temario es el mismo no hace falta hilar especialmente fino para darse cuenta de que la asignatura está mucho más enfocada en el trabajo personal que en modalidad ordinaria. 
%Esta última característica me da pie a hablar del segundo factor: la historia personal del alumnado. 
%La experiencia nos dice que, con frecuencia, el perfil {\footnotesize CIDEAD} ha perdido la inercia del trabajo de estudiante.

 Para acometer el estudio de una asignatura como {\footnotesize FÍSICA} es fundamental recuperar las rutinas necesarias para afrontar los retos que el aparato matemático, probablemente poco lubricado, y el natural ritmo de la asignatura imponen. 

Es difícil establecer un intervalo de horas de trabajo semanal: depende de la situación y recorrido vital de cada cual. Sin embargo no me parece exagerado recomendar la \hl{asistencia presencial} (telemática, \textbf{si no queda más remedio}) a las tutor\'ias y añadir una dedicaci\'on semanal de \hl{no menos de tres horas} para realizar problemas y que comiencen a surgir dudas.


Te recomiendo encarecidamente que distribuyas cuanto antes tu horario semanal de manera conservadora\footnote{Es mejor ir añadiendo minutos poco a poco que \textit{venirse arriba} y hacerte un \textit{planning} en el que la {\scriptsize FÍSICA} ocupe el noventa por ciento de tu tiempo libre. Y menos frustrante cuando la realidad se pasa por casa...} para practicar con regularidad, enfrentar los problemas que se planteen y trabajar los ejercicios entregables para su entrega.


Dicho esto, el \hl{planteamiento general de la asignatura} estar\'a basado en el siguiente esquema de enseñanza-aprendizaje:
\vspace{2ex}

\begin{center}
	{\small
		\sotagboxtres[drop shadow]{\textbf{Tutoría}:\vspace{0pt}	
		 guía y dudas}
		 \vspace{-5ex}
		 
		$$\downarrow$$
		 \vspace{-5ex}
		 
		\sotagboxtres[drop shadow]{\textbf{Trabajo personal}: \\estudio y ejercicios}
		 \vspace{-5ex}
		 
		$$\downarrow$$
		 \vspace{-5ex}
		 
		\sotagboxtres[drop shadow]{\textbf{Atención telemática:} planteamiento de dudas}
	}	
\end{center}

\subsection{Espíritu de la asignatura}
\textbf{i)} La física, como disciplina, es mucho más que calcular numeritos: se trata de una forma de entender los procesos complejos que nos rodean a través de la modelización; \textbf{ii)} la asignatura de {\footnotesize FÍSICA} da un peso importante al trabajo continuo.



\section{Temporización}
Los \hl{contenidos del curso} se desglosan detalladamente en la normativa: debe ser tu base.  Una guía sinóptica de la distribución de las $23$ sesiones de $1.5\,h$ es la siguiente:
\vspace{1ex}

{\footnotesize
	\sotagbox[drop shadow]{1. Ondas en $1+1$ dimensiones}
			\begin{itemize}
				\item Oscilador armónico en {\scriptsize 1D}
	 			\item Ondas armónicas transversales en una cuerda	
	 			\item Ondas de sonido como propagación de un frente esférico 				
	 		\end{itemize}
	 		\vspace{-2ex}
	
	 		{\color{darkgray}{\hspace{1cm}\fbox{{\scriptsize$4$ sesiones ($6\;h$)}}}}

	 		
	 	\sotagbox[drop shadow]{2. Óptica física y geométrica} 
	 		\begin{itemize}
	 			\item Caracterización ($\lambda$, $\nu$, $c$), espectro {\scriptsize EM} e índice de refracción
	 			\item El proceso de transmisión-reflexión. Ecuaciones de \textit{Snell}
	 			\item Óptica geométrica \textit{paraxial}: lentes delgadas y aplicaciones
	 		\end{itemize}
	 		
	 \vspace{-2ex}
	
	 		{\color{darkgray}{\hspace{1cm}\fbox{{\scriptsize$4$ sesiones ($6\;h$)}}}}
	 		 		

	 	 \sotagbox[drop shadow]{3. Campo gravitatorio}
	 		\begin{itemize}
	 			\item Masas, intensidad y potencial.
	 			\item Interacción: energía y trabajo. El problema de \textit{Kepler} circular
	 		\end{itemize}
	 \vspace{-2ex}
	
	 		{\color{darkgray}{\hspace{1cm}\fbox{{\scriptsize$5$ sesiones ($7.5\;h$)}}}}
	 		
 
	 	\sotagbox[drop shadow]{4. Campo electrostático}
	 		\begin{itemize}
	 			\item Cargas eléctricas, intensidad y potencial
	 			\item Interacción: energía y trabajo
	 			\item Conductores y dieléctricos. Ley de \textit{Gau\ss}
	 		\end{itemize}
	 	\vspace{-2ex}
	
	 		{\color{darkgray}{\hspace{1cm}\fbox{{\scriptsize$4$ sesiones ($6\;h$)}}}}
	
			\newpage
		\sotagbox[drop shadow]{5. Campo magnético y electrodinámica}
			\begin{itemize}
				\item Campo magnetostático y aplicaciones
				\item Fuerza de \textit{Heaviside-Lorentz} y ley de \textit{Ampère}
				\item Inducción electromagnética: ley de \textit{Faraday-Henry-Lenz}
			\end{itemize}	
			\vspace{-2ex}
	
	 		{\color{darkgray}{\hspace{1cm}\fbox{{\scriptsize$4$ sesiones ($6\;h$)}}}}
	 		 
	 		

		 \sotagbox[drop shadow]{6. Elementos de relatividad especial}
			\begin{itemize}
				\item Momento y energía relativistas
			\end{itemize}	
		
		\sotagbox[drop shadow]{7. Elementos de física cuántica y nuclear} 
			\begin{itemize}
				\item Dualidad onda-partícula: relación de \textit{de Broglie}
				\item Interacción radiación-materia: efecto fotoeléctrico
				\item Descripción de la desintegración radiactiva
			\end{itemize}	


			\vspace{-2ex}
	
	 		{\color{darkgray}{\hspace{1cm}\fbox{{\scriptsize$2$ sesiones ($3\;h$)}}}}
	 	
}	 		


%%		\item Libro de texto propuesto.
%%			\newline \textit{Util como referencia para la teoría y para los problemas.}
%		\item Alonso, Finn. Física: Campos y ondas (Volumen {\small II}). 
%			\newline \textit{Consulta y ampliación. El nivel del libro es aproximadamente el de un primer curso de grado, aunque asumible como ampliación para este curso. Tomar} cum grano salis.
%	\end{itemize}
%

\section{Evaluación y calificación}

A lo largo del curso trabajaremos con dos instrumentos de evaluación: \textit{exámenes} convocados oficialmente y \textit{ejercicios entregables} cuya realización y entrega se proponen, siempre con carácter voluntario.
	
	\subsection{Exámenes}
	
	Se convocará \hl{un examen al final de cada evaluaci\'on}. Es importante tener en cuenta las siguientes precisiones: 
	{\small
\begin{itemize}
	\item La fecha y hora de realización de cada examen de evaluación será fijada con antelación por
Jefatura de Estudios. A falta de ulterior precisión, los entornos temporales pueden consultarse en el \href{https://www.iesjaimeferran.org/OLD/distancia23/calendario_distancia_2023-24.pdf}{\texttt{calendario de la asignatura.}}
	\item Ante incomparecencia \hl{en ning\'un caso se repetir\'a un examen \textit{sin la debida justificaci\'on documental a cargo de un organismo competente}}.
	\item Si durante la realización de cualquier examen parcial o final un alumno es sorprendido copiando o intercambiando cualquier tipo de información con otro estudiante será calificado con un $0$ en dicha prueba.
	\item Debido a las características de la modalidad el examen emplazado al final de la tercera evaluación hará las veces de examen de recuperación para todos aquellos estudiantes que tengan alguna evaluación suspensa. De este modo la causística hará que coexistan hasta cuatro modelos de examen el mismo día, \textit{vid.}: {\footnotesize I+II+III, I+III, II+III, III}.
	\item Si tienes asignaturas de primer y segundo curso  y vas a presentarte a la convocatoria ordinaria de {\scriptsize EvAU} deberás examinarte de toda la asignatura en algún momento que se fijará entre abril y mayo.
\end{itemize}
}

	\subsection{Ejercicios entregables}
	
	
	En cada evaluaci\'on se propondr\'a la \hl{entrega voluntaria} de ejercicios computable para la calificaci\'on. La realización de los ejercicios podrá hacerse a través de un documento original elaborado por el estudiante o a través de formularios en línea pertinentemente habilitados.
\vspace{1ex}

Son pertinentes las siguientes aclaraciones:
  \begin{itemize}
  {\small
    \item Tanto la publicación de los ejercicios y sus plazos como su entrega se realizarán \textit{exclusivamente} a través de \textit{classroom}. 
    \item 	En caso de requerirse la entrega de documentos originales, la calificación se atendrá a los criterios  de las pruebas escritas.
    \item Los ejercicios correspondientes al tercer trimestre serán tenidos en cuenta en la calificación final de forma alícuota respecto al peso del curso y siempre del  modo ventajoso.
    \item En virtud del espíritu del curso y por a su peso en la evaluación  \hl{se calificar\'a con una nota de $0$ cualquier documento...}: entregado fuera de plazo; cuyo informe de originalidad (en caso de que la entrega consista en un documento elaborado por el estudiante,) \hl{muestre coincidencias claras con otros materiales} alojados en páginas web,  documentos escritos de internet o materiales entregados por otros estudiantes;  \hl{que no se entregue} (en caso de que la entrega consista en un documento original) \hl{en forma de un solo archivo \textit{pdf}, perfectamente legible y bien orientado}.
    }
  \end{itemize}
\vspace{2ex}

	Todos los años me encuentro con gente que saca notazas en los ejercicios voluntarios y notas muy discrepantes en las pruebas presenciales. Además de quedar lamentablemente mal y demostrar una madurez académica por debajo del nivel de lo ridículo, no os va a servir de mucho. Os recomiendo pensar en vuestro orgullo y honor.
	
	
	\subsection{Criterios de calificación específicos}

Tanto para los exámenes como para los ejercicios entregables serán de aplicación los siguientes criterios específicos: 

	{\small
\begin{itemize}
	\item La respuesta a cada ejercicio será calificada con la puntuación máxima cuando la solución esté correctamente planteada, el desarrollo bien justificado y al final se obtenga la solución correcta. 
	\item Se calificarán de cero a diez, en múltiplos de $0.25$.
	\item Se valorará positivamente la realización de esquemas, diagramas o dibujos, y \hl{ser\'a condici\'on indispensable el razonamiento suficientemente detallado de los diferentes pasos}. Asimismo se tendrán en cuenta tanto la destreza en la obtención de resultados numéricos como el proceso seguido en la resolución de las mismas siendo necesaria \hl{la identificaci\'on de los principios y leyes involucrados}.
	\item Será necesario presentar los resultados con las \hl{unidades adecuadas} así como respetar la \hl{naturaleza vectorial o escalar} de las magnitudes con las que se opera. Cada falta o incorrección en su asignación se penalizará con $0.25$ puntos hasta un máximo de $2$ puntos en el examen.
	\item Los \hl{errores ortogr\'aficos ser\'an tenidos en cuenta} para la calificación según el criterio siguiente: hasta dos errores aislados no deben penalizarse; diferentes y sucesivos errores ortográficos se penalizarán con un descuento de $0.25$ cada uno hasta un máximo de $2$ puntos a descontar sin perjuicio de la penalización aplicada en el apartado anterior.
	\item Se ponderará específicamente la \hl{capacidad expresiva y la correcci\'on idiom\'atica} de los estudiantes. Para ello se tendrá en cuenta:   la propiedad del vocabulario y la corrección sintáctica; la puntuación apropiada; la adecuada presentación.
\end{itemize}
}
\newpage

\subsection{Cálculo de calificaciones}
La nota se obtendrá por \hl{la m\'as ventajosa} de estas vías:
\vspace{2ex}

\hspace{.5cm}{\small
\begin{minipage}{.5\textwidth}
	\sotagboxtwo[drop shadow]{
	\center{
	\textbf{Opción A}: \vspace{1ex}
	
	{\footnotesize
	 Nota del examen de la evaluación ($75\%$) y la entrega de ejercicios ($25\%$)}}}
\end{minipage}\hfill
\begin{minipage}{.5\textwidth}
	\sotagboxtwo[drop shadow]{
	\center{
	\textbf{Opción B}:\vspace{1ex}
	
	{\footnotesize
	 Nota del examen\\ de evaluación}}\vspace{.5cm}
	 
	 }
\end{minipage}
}

\vspace{3ex}

\begin{itemize}
\item Para poder acceder al $25\%$ procedente de entregables será necesario presentarse al examen correspondiente.
\end{itemize}


\vspace{1ex}

\subsubsection[Nota final]{Cálculo de la nota final}

La nota final de la asignatura en convocatoria ordinaria se calcula considerando la siguiente casuística:

\vspace{1ex}

	{\small
	
	\textbf{\hl{Si has aprobado las tres evaluaciones:}}
	\vspace{-1ex}
	
	\begin{itemize}
		\item[\textit{a)}] La nota de final de curso será la media aritm\'etica de las notas de las tres evaluaciones. 
	\end{itemize}
	
	\textbf{\hl{Si has suspendido al menos una evaluaci\'on:}}
	\vspace{-1ex}
	\begin{itemize}
		\item[\textit{b)}] La nota final del curso será proporcional a los contenidos de la prueba que enfrentes en el examen de la tercera evaluación. De esta forma:
			\begin{itemize}
			
				\item \footnotesize{te presentas a I+III, II+III}
				\[
				\text{\small nota final}=\frac{2}{3}\,
					\begin{cases}
						\text{nota {\footnotesize(I+III)}}\\
						\text{nota {\footnotesize(II+III)}}
					\end{cases}
									+\frac{1}{3}\,
					\begin{cases}
						\text{ nota{\footnotesize(II)}}\\
						\text{ nota{\footnotesize(I)}},
					\end{cases}		
				\]							
				\item \footnotesize{te presentas a I+II+III}
				\[
				\text{\small nota final}=\text{\small nota media\footnotesize{(I+II+III)}}.
				\]
			
			\end{itemize}
	
	\end{itemize}
 }
 
  El alumno estará aprobado si la nota final, no redondeada, calculada según \textit{a)} o  \textit{b)} es igual o superior a $5$. En caso de no aprobar en convocatoria ordinaria el estudiante tendrá la oportunidad de presentarse a la convocatoria extraordinaria en un examen global.

\newpage
\thispagestyle{empty}
\null
\newpage
\small{
\begin{versionhistory}
 	 	\vhEntry{1.0}{08.09.22}{R. Torres}{Liberado para su publicación.}
  	 \vhEntry{1.2}{08.09.23}{R. Torres}{Eliminado el {\footnotesize QR} para matriculación, suprimida cualquier referencia a {\footnotesize CIDEAD}. Añadidos comentarios respecto a cuestionarios de evaluación \textit{online} y sobre exámenes anticipados.}
  	  \vhEntry{1.3}{05.09.24}{R. Torres}{Incluidos comentarios explícitos acerca de las matrículas parciales y la gestión de las tutorías.}
  	  \vhEntry{1.4}{15.09.25}{R. Torres}{Cambios en la calificación de los ejercicios entregables.}
  	 \vhEntry{1.5}{07.09.26}{R. Torres}{Mención explícita al carácter de guía de la tutoría, así como de que no es una clase de teoría.}
\end{versionhistory}
	}
	
	
	
\end{document}

